% !TeX root = ../papers/nc-p3.tex
\section{不変直線から相補逆行列窓への全持上げと二次障害の個数}
\label{sec:invariant-line-filtered-lifts}
可逆二平面分解の任意の窓の四テンソルを $w_2$ とする。
二文字ずつの空間を $N_i,T_i$ とし、$N$ の文字数を重みとする。
$w_2$ は重み2であり、$N_i\otimes N_{i+1}$ 内の関係を $c_i$ と書く。
四位置の文字列の長さは常に4であり、以下の添字は標準四移動で読む。

\begin{theorem}\label{thm:invariant-line-filtered-lifts}
$w_2$ に指定した二平面の重みで特殊化する全窓は、重み3の64係数 $u$ と
重み4の16係数 $v$ から $w_2+u+v$ として得られる。許容条件は次の通りである。
\begin{enumerate}
\item $u$ には48行の斉一次方程式があり、その解空間の次元は16以上24以下である。
\item $c_0,c_1,c_2$ を交換子にする基底を順に取り、残る $c_3$ の対称部分を $q$ とする。
$q=0$（四つの射影的な同一視の合成が恒等）のとき、許される $u$ は
この一次解空間上の五成分の斉二次障害の零点で、各 $u$ の上の $v$ は三次元のアフィン空間である。
$q\ne0$ なら二成分の斉二次障害となり、各零点の上の $v$ は一意である。
成分数は標的の次元を表し、二次式の独立性や高さを要求しない。
\item 全解は全次数で標準Koszul左右AS正則で、指定直線上で前後とも評価階数3である。
指定直線を保つ添字同型は下記の線型群とせん断の軌道である。
従ってこの構成と $w_2$ の相補逆行列の正規形で、直線外に分岐を許した
定理\ref{thm:four-shift-line-pbw}の正則枝全体を被覆する。
\end{enumerate}
\end{theorem}
\begin{proof}[二切断の消去]
切断 $i,i+1\mid i+2,i+3$ を $i=0,1$ で取る。各二文字位置を $TT$, 混合, $NN$ の順に書くと行列は
\[
 \begin{pmatrix}0&0&A\\0&B&C\\D&E&F\end{pmatrix}.
\]
$A,D$ が階数1、$B$ が階数4で、これらは $w_2$ の係数である。
$C,E$ は $u$、$F$ は $v$ だけで決まる。$A$ の $NN$ 側の直線と
$D$ の $NN$ 側の直線がちょうど $k c_{i+2},k c_i$ である。
この両直線で割った $C,E,F$ の像を $C',E',F'$ とする。行・列のピボット二つを外すと
残るこの行列の階数が4である条件は、$BJB=B$ なる一つの一般逆 $J$ を用いて
\[
 \operatorname{im} C'\subset\operatorname{im} B,\qquad
 \operatorname{row} E'\subset\operatorname{row} B,\qquad F'=E'JC'
\]
である。前二条件の下で右辺は一般逆の選択に依存しない。
各切断で一次条件は $4\times3$ 個ずつ二組、二次の右辺は $3\times3$ 個である。

一次条件の余核次元は高々8である。実際、接線文字が位置 $j$ だけにある十六空間を $U_j$ とする。
辺 $j,j+1$ の一次条件は $U_j\oplus U_{j+1}$ から十二空間への全射で、各直和因子への制限も全射である。
これは混合関係四平面が両混合半分に同型に射影し、反対辺では $NN$ を $kc_{j+2}$ で割るからである。
四つの条件の線型従属では、どれか一辺の係数で残り三辺の係数が決まる。
実際隣り合う条件の係数は $U_j^*$ 内で互いに打ち消し、各々への双対写像は単射である。
その共通像は $U_j^*$ の中で三法線文字の相続く二辺の $c$ をそれぞれ零化する。
これら二辺を交換子に揃えると、三つの法線文字で完全に対称な三次式と接線文字の積だから次元8である。
これで四十八行の階数は40以上48以下である。

次に $v$ の二つの射影の合計を $L:k^{16}\to k^9\oplus k^9$ とする。
各射影の双対像は、互いに反対の二辺でそれぞれ $c$ を零化する四テンソルの九空間である。
二つの像の交わりは四辺全てで $c$ を零化する。最初の三辺で交換子を零化すれば
通常の $\operatorname{Sym}^4(k^2)^*$ となり、残る縮約は $q$ による二次の縮約である。
$q\ne0$ ならその双対は二次式への $q$ 倍で単射だから、縮約の階数は3である。
従って交わりの次元は $q=0$ なら5、他は2、$L$ の階数はそれぞれ13と16である。
二つの $E'JC'$ の組の $\operatorname{coker}L$ での類が主張の明示的な二次障害となる。
\end{proof}
\begin{proof}[全次数性と同型]
上の式から各切断は階数6である。重みの一径数で $u,v$ にそれぞれ $t,t^2$ を掛けると、
全ての $t\ne0$ で同じ窓、$t=0$ で $w_2$ が得られる。
命題\ref{prop:unrestricted-good-reduction}から非零ファイバーも標準Koszul左右AS正則である。
直線上の評価を調べる。関係空間の重みによる初期空間は $w_2$ の関係空間である。
実際二つの切断はいずれも階数6一定なので、像の極限は特殊点の像である。
従って接線二文字への射影はその行列式であり、混合四関係の持上げの評価は
任意の指定直線上の点で二つの法線文字を張る。左右とも階数3である。
逆に指定不変直線を持つ初めの正則窓では定理\ref{thm:line-polystable-regular-limit}の
最初の特殊化がちょうどこの $w_2$ を与え、係数は以上の全条件を満たす。

$w_2$ の射影類を保つ $\prod_i(\mathrm{GL}(N_i)\times\mathrm{GL}(T_i))$ と
$h_i:T_i\to N_i$ のせん断で全ての指定付き同型が得られる。
先の群では $w_2$ の代表を揃えるため全テンソルの共通倍率を正規化する。
せん断の作用は各 $h_i$ をその位置だけに作用させて
\[
 u\mapsto u+\textstyle\sum_i h_iw_2,\qquad
 v\mapsto v+\textstyle\sum_i h_i u+\sum_{i<j}h_i h_j w_2
\]
である。これは指定直線を保つ全基底群のブロック三角分解から必要十分で、
余分な変換は生じない。印を忘れるときは第二の窓の指定可能な四移動不変直線を動かして比較する。
一次元グラフならこの選択は成分数以下である。一移動・二移動にも同じ比較を適用できる。
\end{proof}
